\input{../../resources/lesson-head.tex}

\subsection*{Week 4 at a glance}

\subsubsection*{We will be learning and practicing to:}
%evaluating compound propositions
%truth table definitions
%consistency
%translating
%logical equivalence via laws
%logical equivalence via truth tables
%variants of conditionals
\begin{itemize}

\item Translate between different representations to illustrate a concept.
\begin{itemize}
   \item Translating between symbolic and English versions of statements using precise mathematical language
   \item Translating between truth tables (tables of values) and compound propositions
\end{itemize}

\item Use precise notation to encode meaning and present arguments concisely and clearly
\begin{itemize}
    \item Listing the truth tables of atomic boolean functions (and, or, xor, not, if, iff)
    \item Defining functions, predicates, and binary relations using multiple representations
\end{itemize}

\item Know, select and apply appropriate computing knowledge and problem-solving techniques. Reason about computation and systems. Use mathematical techniques to solve problems. Determine appropriate conceptual tools to apply to new situations. Know when tools do not apply and try different approaches. Critically analyze and evaluate candidate solutions.
\begin{itemize}
    \item Evaluating compound propositions
    \item Judging logical equivalence of compound propositions using symbolic manipulation with known equivalences, including DeMorgan's Law
    \item Writing the converse, contrapositive, and inverse of a given conditional statement
    \item Determining what evidence is required to establish that a quantified statement is true or false
    \item Evaluating quantified statements about finite and infinite domains
\end{itemize}

\end{itemize}

\vfill


\subsubsection*{TODO:}
\begin{list}
   {\itemsep2pt}
    \item Schedule your Test 1 Attempt 1, Test 2 Attempt 1 times at PrairieTest (http://us.prairietest.com)   
    \begin{itemize}
    \item You must present a physical university-issued or government-issued ID. Copies, photos, or digital IDs are not accepted, and students without verifiable ID will not be permitted to test. 
    \item Pockets must be empty upon entry. Students may bring only their ID and a pen or pencil into the testing center. 
    \end{itemize}

    \item Review quizzes based on Week 2 and Week 3 class material (due Monday 01/26/2026)
    \item Homework 2 due on Gradescope https://www.gradescope.com/ (Thursday 01/29/2026)
   \item Start reviewing for Test 1: {\bf Test1 Attempt 1 is next week}. The test covers material in Weeks 1 through Week 5, corresponding to RQ1, RQ2, RQ3, and RQ4. To study for the exam, we recommend reviewing class notes (e.g. annotations linked on the class website, supplementary videos from the class website) and working multiple variants of review quiz questions. You could also find that reviewing homework (and its posted sample solutions)and extra examples (in lecture notes and discussion examples) and getting feedback (office hours and Piazza) is helpful.
\end{list}

\newpage

\section*{Week 4 Monday: Logical Equivalence and Conditionals}
\input{../activity-snippets/logical-equivalence-identities-andornotxor.tex}
\newpage
\input{../activity-snippets/logical-operators-full-truth-table.tex}
\input{../activity-snippets/hypothesis-conclusion.tex}
\input{../activity-snippets/converse-inverse-contrapositive.tex}
\vfill
\input{../activity-snippets/compound-propositions-recursive-definition.tex}
\input{../activity-snippets/compound-propositions-precedence.tex}
\newpage
%! app: TODOapp
%! outcome: logical equivalence via laws, logical equivalence via truth tables

{\bf (Some) logical equivalences}

\begin{tabular}{llp{3in}}
$p \to q \equiv \lnot p \lor q$ & & \\
&& \\
&& \\
$p \to q \equiv \lnot q \to \lnot p$ & &{\bf Contrapositive} \\
&& \\
&& \\
$\lnot (p \to q) \equiv p\land \lnot q$  & &\\
&& \\
&& \\
$\lnot( p \leftrightarrow q) \equiv p \oplus q$ && \\
&& \\
&& \\
$p \leftrightarrow q \equiv q \leftrightarrow p$ &&\\
&& \\
\end{tabular}

\vfill

{\it Extra examples}:

$p \leftrightarrow q$ is not logically equivalent to $p \land q$ because \underline{\phantom{\hspace{4in}}} 

$p \to q$ is not logically equivalent to $q \to p$ because \underline{\phantom{\hspace{4in}}} 
\vfill

\newpage
\input{../activity-snippets/logical-operators-english-synonyms.tex}
\newpage


\section*{Week 4 Wednesday: Translation, Predicates, and Quantifiers}
\input{../activity-snippets/compound-propositions-translation.tex}
\newpage
\input{../activity-snippets/consistency-def.tex}
\input{../activity-snippets/consistency-example.tex}
\newpage
\input{../activity-snippets/predicate-definition.tex}
\input{../activity-snippets/predicate-examples-finite-domain.tex}
\vfill
\input{../activity-snippets/predicate-truth-set-definition.tex}
\input{../activity-snippets/predicate-truth-set-example.tex}
\newpage
\input{../activity-snippets/quantification-definition.tex}
\vfill
\input{../activity-snippets/quantification-logical-equivalence.tex}
\vfill
\vfill
%! app: TODOapp
%! outcome: TODOoutcome

An element for which $P(x) = F$ is called a {\bf counterexample} of $\forall x P(x)$. When there is a counterexample in the domain of predicate $P$,  $\forall x P(x)$ is false.


An element for which $P(x) = T$ is called a {\bf witness} of $\exists x P(x)$. When there is a witness in the domain of predicate $P$, then $\exists x P(x)$ is true.

\vfill
\input{../activity-snippets/quantification-examples-finite-domain.tex}
\vfill
\newpage

\section*{Week 4 Friday: Evaluating Nested Quantifiers}
\input{../activity-snippets/rna-rnalen-basecount-definitions.tex}
%! app: TODOapp
%! outcome: evidence for quantified statements

{\bf Example predicates on $S$, the set of RNA strands (an infinite set)}


$H: S \to \{T, F\}$ where $H(s) = T$ for all $s$.

Truth set of $H$ is \underline{\phantom{$S$\hspace{1in}}}

\vfill

$L_{\A}: S \to \{T, F\}$  defined recursively by: 

~~Basis step: $L_{\A}(\A) = T$, $L_{\A}(\C) = L_{\A}(\G) = L_{\A}(\U) = F$

~~Recursive step: If $s \in S$ and $b \in B$, then $L_{\A}(sb) = L_{\A}(s)$.

Example where $L_{\A}$ evaluates to $T$ is \underline{\phantom{$\A\C\G$~\hspace{1in}}}

\vfill

Example where $L_{\A}$ evaluates to $F$ is \underline{\phantom{$\U\A\C\U$\hspace{1in}}}

\vfill
\newpage
\input{../activity-snippets/predicate-notation.tex}
\vfill
\input{../activity-snippets/predicates-example-rnalen-basecount.tex}
\newpage
\input{../activity-snippets/predicates-projecting-example-rna-basecount.tex}
\vfill
\input{../activity-snippets/nested-quantifiers.tex}
\newpage
\input{../activity-snippets/alternating-quantifiers-strategies-rna-examples.tex}
\newpage

\begin{comment}
\subsection*{Review Quiz}
\begin{enumerate}
\item Logical equivalence
 \input{../activity-snippets/quiz-truth-values-conditional.tex}
\item Translating propositional logic
    \begin{enumerate}
    \item \input{../activity-snippets/quiz-compound-propositions-translation.tex}
    \newpage
    \item \input{../activity-snippets/quiz-consistency.tex}
    \end{enumerate}
\newpage
\item Evaluating predicates
    \begin{enumerate}
    \item \hspace{1in}\\ \input{../activity-snippets/quiz-predicates-finite-domain.tex}
    \item \hspace{1in}\\ \input{../activity-snippets/quiz-predicates.tex}
    \item \hspace{1in}\\ %! app: Bioinformatics
%! outcome: recursive definitions, function and relation definitions

For this question, we will use the following predicate.

$L_{\A}$ with domain $S$ is defined recursively by: 
\begin{itemize}
\item[]Basis step: $L_{\A}(\A) = T$, $L_{\A}(\C) = L_{\A}(\G) = L_{\A}(\U) = F$
\item[]Recursive step: If $s \in S$ and $b \in B$, then $L_{\A}(sb) = L_{\A}(s)$
\end{itemize}

Which of the following is true? (Select  all and only that apply.)
 \begin{enumerate}
    \item $L_{\A} ( \A\A)$
    \item $L_{\A} ( \A\C)$
    \item $L_{\A} ( \A\G)$
    \item $L_{\A} ( \A\U)$
    \item $L_{\A} ( \C\A)$
    \item $L_{\A} ( \C\C)$
    \item $L_{\A} ( \C\G)$
    \item $L_{\A} ( \C\U)$
 \end{enumerate}    

    \end{enumerate}
\newpage
\item Evaluating nested predicates
    \begin{enumerate}
    \item \hspace{1in}\\\input{../activity-snippets/quiz-predicates-alternating-quantifiers-rnalen.tex}
    \item \hspace{1in}\\\input{../activity-snippets/quiz-predicates-alternating-quantifiers-basecount.tex}
    \end{enumerate}
\end{enumerate}
\end{comment}
\end{document}
